\section{Failure of the classical total-mass calibration for relative extremals}
\paragraph{Status and scope.}
This is an original-category obstruction to importing the classical
relative-extremal-function proof of Borel capacity regularity unchanged.
It is not a counterexample to B1. The background is smooth and strongly
positive, $q=1$, and the original dimension is any $n\ge2$. The exact
capacity calculation uses every original normalized classical psh tester,
not only toric testers. No equality of an actual envelope with its toric
trace is assumed, and no nonlinear pushforward identity is used.

Fix the prescribed finite $K/\mathbf Q_p$ and its absolute value,
$\lambda=-\log|\pi|>0$, the original torus $X=(\mathbf G_m^n)^{\rm an}$,
$t_i=\log|T_i|$, and
\[
 V=\{|t_i|<\lambda\ (1\le i\le n)\},\quad
 \beta=\sum_{i=1}^n d't_i\wedge d''t_i,\quad
 S=\beta^{n-1}\wedge[X],\quad q=n-(n-1)=1.
\]
Use normalized moments $x_i=t_i/\lambda$ and write
$\tau=(x_1,\ldots,x_n):V\to Q=(-1,1)^n$. Let
$\sigma:Q\to V$ be the joint original Gauss section and set
\[
 \gamma=(n-1)!\lambda^{n-2},\qquad A_n=n!2^n\lambda^{n-2}.
\]
For $0<r<1$, define
\[
 K_r=\tau^{-1}([-r,r]^n),\quad
 a_r=\max\left\{{\|x\|_\infty-r\over1-r},0\right\},\quad h_r=a_r-1.
 \tag{EC1}
\]
The compact original analytic domain $K_r$ has nonempty ordinary interior.
The functions $a_r$ and $h_r$ are actual continuous classical psh functions
formed from finite real maxima of original harmonic affine potentials on
the same ordinary representative. In particular $0\le a_r<1$ on $V$.

\begin{lemma}[A local linear trace identity for the full bounded tester class]
Let $a$ be any original bounded classical psh function on $V$. The bounded
Borel function $a\circ\sigma$ has a unique continuous convex representative
$g_a$ modulo Lebesgue null sets. As Radon measures on $Q$,
\[
 \tau_*(Da\wedge S)=\gamma\Delta g_a.
 \tag{EC2}
\]
Pointwise inequalities for actual functions pass to their canonical
convex traces. No assertion about arbitrary completed Borel representatives
is involved here.
\end{lemma}
\begin{proof}
For $\varphi\in C_c^\infty(Q)$, $\varphi\circ\tau$ is an original compactly
supported smooth CLD scalar test: the moment map is proper over compact
subsets of $Q$. The current recursion, compact Stokes, and the frozen full
original top-form Gauss calibration yield
\[
 \begin{split}
 \int_V(\varphi\circ\tau)\,d(Da\wedge S)
 &=\langle aS,D(\varphi\circ\tau)\rangle\\
 &=(n-1)!\lambda^{n-2}\int_Q(a\circ\sigma)\Delta\varphi\,dx.
 \end{split}
 \tag{EC3}
\]
The coefficient follows from
$D\varphi\wedge\beta^{n-1}=(n-1)!\Delta\varphi\,d\mathrm{vol}$
in unnormalized moment coordinates: the rescaling contributes
$\lambda^{-2}$ from differentiation and $\lambda^n$ from volume.

For completeness, convexity of the canonical trace need not be imported
from a completed zero-boundary theorem. Fix a unit vector $v\in\mathbf R^n$
and an orthonormal basis $w_2,\ldots,w_n$ of $v^\perp$. The constant
strongly positive $(n-1,n-1)$ form
\[
 \Theta_v=\bigwedge_{j=2}^n
 d'\langle w_j,t\rangle\wedge d''\langle w_j,t\rangle
\]
is separately closed. The positive bounded BT measure
$Da\wedge\Theta_v\wedge[X]$, tested against nonnegative
$\varphi\circ\tau$, gives by the same Stokes and top-form computation
\[
 0\le\lambda^{n-2}\int_Q(a\circ\sigma)\,\partial_{vv}\varphi\,dx.
\]
Thus every distributional second directional derivative of
$a\circ\sigma$ is nonnegative. Local mollifications have positive
semidefinite Hessian and are convex. They remain uniformly bounded on
protected compact boxes; the elementary interior slope bound for bounded
convex functions makes them locally equicontinuous. Subsequential locally
uniform limits produce a continuous convex representative $g_a$; convergence
in distributions identifies it with $a\circ\sigma$ almost everywhere.
Any two such representatives agree everywhere by continuity. Bounds and
order inequalities pass first almost everywhere and then everywhere.
This also proves the asserted canonical nature of the trace without
asserting equality at every individual Gauss point.

Equation (EC3) now identifies the two positive Radon measures in (EC2).
Only local boundedness, current recursion, positivity, and the full
original top-form calibration have been used. The energy and zero-boundary
assumptions imposed in R17 FP for its later conclusions are not used.
This is a linear degree-one identity, not an arbitrary higher-degree
Monge--Ampere pushforward assertion.
\end{proof}

\begin{proposition}[Exact capacity over all original normalized testers]
For every $0<r<1$,
\[
 \boxed{\quad C_S(K_r;V)=C_S^*(K_r;V)
             =A_n\,{r^{n-1}\over1-r}.\quad}
 \tag{EC4}
\]
\end{proposition}
\begin{proof}
Let $a$ be any original tester with $0\le a\le1$. Then
$0\le g_a\le1$. Fix all coordinates other than $x_i$. The resulting
one-variable convex function $f$ has
\[
 f''([-r,r])=f'_+(r)-f'_-(-r)\le {2\over1-r}.
 \tag{EC4a}
\]
Indeed the right derivative is at most
$(1-f(r))/(1-r)\le1/(1-r)$ by comparison with secants approaching $1$;
the left derivative is at least
$(f(-r)-1)/(1-r)\ge-1/(1-r)$ by secants approaching $-1$.
Distributional slicing, obtained by Fubini from the definition of
$\partial_{ii}g_a$, integrates (EC4a) over the other $n-1$ coordinates
and gives
\[
 \partial_{ii}g_a([-r,r]^n)\le {2(2r)^{n-1}\over1-r}.
\]
Summing over $i$ and multiplying by $\gamma$ in (EC2) gives the upper
bound in (EC4) for every original tester. This is a slice computation,
not an assertion of an $n$-dimensional atom on a coordinate hyperplane.

For the explicit $a_r$, fix the other $n-1$ coordinates in the inner
cube. In the remaining coordinate the function is
$\max\{(|x_i|-r)/(1-r),0\}$, so its derivative measure on $[-r,r]$
has mass $2/(1-r)$, including its two endpoint atoms. Integrating these
slice measures gives exactly $2(2r)^{n-1}/(1-r)$ in each coordinate.
This also accounts for all edges and corners, without any extra mass.
Therefore
\[
 \Delta a_r([-r,r]^n)={2n(2r)^{n-1}\over1-r},
\]
so $a_r$ attains the full raw upper bound. For $r<s<1$, the open set
$\tau^{-1}((-s,s)^n)$ contains $K_r$ and its capacity is at most the just
computed value $C_S(K_s;V)$. Let $s\downarrow r$ and use
$C_S\le C_S^*$ to obtain the outer equality. This special exact
calculation does not revive universal compact raw/outer equality.
\end{proof}

\begin{proposition}[Obstruction to total-mass extremal calibration]
Every bounded actual classical psh function $u$ satisfying
\[
 h_r\le u\le0\text{ on }V,\qquad
 u\le-1\text{ on }\operatorname{Int}K_r
 \tag{EC5}
\]
has canonical convex trace $g_u=h_r\circ\sigma$ and
\[
 \begin{split}
 \int_V Du\wedge S&={A_n\over1-r},\\
 \int_{V\setminus K_r}Du\wedge S
   &=A_n{1-r^{n-1}\over1-r}>0,\\
 {\displaystyle\int_V Du\wedge S\over C_S^*(K_r;V)}
   &=r^{1-n}.
 \end{split}
 \tag{EC6}
\]
In particular no such assignment of actual extremals can identify their
total relative source masses with the original capacities, even up to a
fixed multiplicative upper constant as the obstacle $K_r$ varies within
this one fixed original domain and smooth background.
\end{proposition}
\begin{proof}
The order-preserving trace lemma and continuity show $g_u=-1$ on
$[-r,r]^n$ and $g_u\le0$ on $Q$. If $R=\|x\|_\infty>r$, intersect the
ray through $x$ with the cubes of radius $r$ and $s$, where $R<s<1$.
Convexity between these two points, followed by $s\uparrow1$, gives
\[
 g_u(x)\le {R-1\over1-r}.
\]
The lower bound $u\ge h_r$ gives the reverse inequality. This proves
the trace identity without asserting equality off the Gauss section.

For $r<s<1$, flux through the cube of radius $s$ gives
\[
 \Delta g_u([-s,s]^n)={2n(2s)^{n-1}\over1-r}.
\]
Equivalently, the same slicing computation gives derivative jump
$2/(1-r)$ in each coordinate for almost every fixed choice of the other
coordinates in $[-s,s]^{n-1}$; integration yields the displayed formula
without invoking a boundary regularity theorem. Letting $s\uparrow1$ uses
only monotone convergence for this single Radon measure, not continuity
of an outer capacity. The inner mass is the value already computed in
(EC4). Multiplying by $\gamma$ and subtracting proves (EC6). For
$n\ge2$, $r^{1-n}\to\infty$ as $r\downarrow0$, proving the last assertion.
\end{proof}

\paragraph{Application to a putative relative extremal function.}
Let
\[
 H=\sup\{v:v\text{ is original classical psh},\ v\le0\text{ on }V,
                                      \ v\le-1\text{ on }K_r\}.
\]
Replacing any competitor by $\max\{v,-1\}$ increases it and preserves
the constraints, so the envelope can be taken over bounded competitors.
The explicit $h_r$ is a competitor; hence $h_r\le H\le0$ and $H=-1$ on
$K_r$. Its ordinary upper semicontinuous regularization $H^*$ satisfies
$h_r\le H^*\le0$ and $H^*=-1$ on $\operatorname{Int}K_r$.
If $H^*$ is an actual classical psh function, (EC6) applies to it.
Consequently its source has positive mass away from the compact obstacle
and its total mass cannot equal (EC4). If psh-envelope existence is not
available, the conjunction of that existence assertion and the classical
total-mass calibration is still refuted. No envelope-existence theorem,
all-fiber envelope identity, or completed point representative is assumed.

\paragraph{Precise implication for B1.}
One cannot import the classical extremal proof simply by replacing
$(dd^c h)^n$ with $(Dh)^q\wedge S$ while keeping the original classical
psh competitor cone. Already for arbitrary original dimension $n\ge2$
and remaining degree $q=1$, the forced convex trace carries relative
source mass outside its obstacle. Enlarging the competitor cone to all
functions subharmonic for the trace operator changes the tester problem
and is not an admissible repair. Original outer Borel compact-inner
regularity remains OPEN; this is a failed analytic calibration, not a
failure of B1 itself.

\paragraph{Internal dependencies and review boundary.}
The proof uses the frozen bounded fixed-background BT recursion and
positivity, max stability of bounded classical psh functions, and the
full original top-form Gauss calibration (R17 FP6, source
8389d2ce81aa5fd1b563767a3c207a06d0b50e661fb06ee978da2d39634c49be.tex,
with its local calculation explicitly extended above). The directional
constant forms give a self-contained convex-trace argument. No energy
completion, abstract weak-compactness inference, general GQ/KC12 descent,
model/conormal restriction, or original raw/outer equality is imported.
